\section{Objet et statut de l'annexe}
\label{annexe-histoires-objet}

Cette annexe fixe le cadre formel utilisé dans le Tome III pour parler
sans ambiguïté des histoires, des observations et des classes de
systèmes qu'elles permettent de distinguer. Elle rassemble les
hypothèses communes aux résultats de séparation et établit les quatre
propriétés auxquelles les chapitres renvoient ensuite :

\begin{itemize}
\item
  le transport d'une équivalence par un préfixe, pour les lectures compatibles ;
\item
  la congruence comportementale obtenue lorsque toutes les
  continuations admissibles sont disponibles ;
\item
  la monotonie conjointe par enrichissement du budget d'histoires et
  raffinement de la résolution ;
\item
  l'indépendance du quotient vis-à-vis de la mesure de complexité
  choisie pour graduer les histoires.
\end{itemize}

Le point de départ est volontairement abstrait. Les épisodes peuvent
être des expositions, des attentes, des regraînes ou des protocoles
composés. Les observations peuvent être binaires, compter certains
événements ou lire une partie plus riche de l'exécution. Les résultats
ci-dessous ne dépendent pas de la réalisation numérique particulière
du modèle, mais seulement de la cohérence de la composition des
histoires et de la détermination des réponses.

\begin{encadre}[title={Principe de relativité observationnelle}]
Aucune classe, aucune profondeur de séparabilité et aucun résumé suffisant
n'est attaché à un système isolément. Ces objets sont toujours relatifs
au couple
\[
(\mathfrak B,\rho),
\]
où \(\mathfrak B\) déclare ce que l'on peut faire vivre au système et
\(\rho\) ce que l'on choisit d'observer.
\end{encadre}

\section{Systèmes, épisodes et histoires}
\label{annexe-histoires-monoide}

Soit \(\mathfrak X\) l'ensemble des états complets ou, plus
généralement, des présentations de systèmes admises dans le Tome III.
On note \(\mathcal E\) l'ensemble des épisodes élémentaires. Conformément
à la décision D-06, cet alphabet comprend les épisodes actifs et les
délais :
\[
\mathcal E=\mathcal E_{\mathrm{act}}\cup\mathcal D,
\]
où \(\mathcal D\) désigne les générateurs d'attente. Un délai est donc
une opération du protocole à part entière et non une absence d'action :
il laisse évoluer les horloges internes et peut modifier la réponse
ultérieure du système.

Une histoire finie est un mot de l'ensemble
\[
\mathfrak M=\mathcal E^{\ast}.
\]
L'histoire vide est notée \(\varepsilon\). La concaténation de deux
histoires \(\mathcal H_1\) et \(\mathcal H_2\) est notée
\(\mathcal H_1\cdot\mathcal H_2\). Le triplet
\[
(\mathfrak M,\cdot,\varepsilon)
\]
est le monoïde libre engendré par \(\mathcal E\).

À chaque épisode \(e\in\mathcal E\) est associée une application de
transition
\[
\Phi_e:\mathfrak X\longrightarrow\mathfrak X.
\]
Pour une histoire
\[
\mathcal H=e_1e_2\cdots e_n,
\]
on définit
\[
\Phi_{\mathcal H}
=
\Phi_{e_n}\circ\cdots\circ\Phi_{e_1},
\qquad
\Phi_{\varepsilon}=\operatorname{Id}_{\mathfrak X}.
\]

La convention de composition adoptée dans l'ensemble du Tome III est
donc
\begin{equation}
\label{eq:annexe-G-composition}
\boxed{
\Phi_{\mathcal H_1\cdot\mathcal H_2}
=
\Phi_{\mathcal H_2}\circ\Phi_{\mathcal H_1}.
}
\end{equation}
Une histoire se lit de gauche à droite dans l'ordre chronologique : le
système vit d'abord \(\mathcal H_1\), puis \(\mathcal H_2\).

Cette convention est importante pour les classes de continuations
introduites plus loin. Elle évite aussi de confondre la composition
temporelle des protocoles avec la convention usuelle de composition des
fonctions.

\begin{encadre}[title={A-H1 --- statut de la regraine}]
Lorsqu'une regraine est appliquée, elle appartient au protocole, donc à
l'histoire, et non à la description de l'état initial. Deux expériences
qui diffèrent par leur regraine correspondent ainsi à deux histoires
distinctes, même lorsqu'elles partent du même état. Cette convention sera
utilisée dans les annexes consacrées à la réduction et à la suffisance de
l'état lent.
\end{encadre}

\section{Protocoles et classes admissibles}
\label{annexe-histoires-admissibilite}

Un protocole expérimental ne donne pas nécessairement accès à tout
\(\mathfrak M\). Il déclare une classe
\[
\mathfrak B\subseteq\mathfrak M
\]
d'histoires admissibles. Dans les applications, \(\mathfrak B\) peut
être fini, paramétré par une borne de durée ou défini par plusieurs
contraintes simultanées.

La seule hypothèse requise pour définir un quotient d'observation est
que chaque histoire de \(\mathfrak B\) possède une exécution et une
lecture bien définies. Des propriétés de fermeture supplémentaires ne
sont introduites que lorsqu'un résultat les utilise effectivement.

\subsection{Trois niveaux de fermeture}
\label{annexe-histoires-fermetures}

On distingue les niveaux suivants.

\paragraph{Classe minimale.}
Une classe minimale est un sous-ensemble déclaré de \(\mathfrak M\),
sans hypothèse de fermeture. Elle suffit pour comparer les réponses sur
un protocole fini, établir la monotonie des quotients et démontrer les
résultats de réduction qui portent sur une sonde ou une famille de
sondes explicitement déclarée.

\paragraph{Classe préfixe-fermée.}
La classe \(\mathfrak B\) est préfixe-fermée si, pour toute
décomposition
\[
\mathcal H=\mathcal H_1\cdot\mathcal H_2\in\mathfrak B,
\]
on a \(\mathcal H_1\in\mathfrak B\). Cette propriété garantit que les
états intermédiaires d'une histoire admissible sont eux-mêmes atteints
par une histoire du protocole.

Les budgets définis comme ensembles inférieurs d'une complexité croissante par préfixe sont automatiquement préfixe-fermés. C'est notamment le cas des contraintes cumulatives usuelles portant sur le nombre d'épisodes, la durée totale ou la dose reçue. Le niveau préfixe-fermé est donc naturel pour ces budgets sans être requis par la définition générale d'un quotient d'observation.

\paragraph{Classe stable par continuation.}
La classe est stable par continuation dans le monoïde considéré si,
pour tous \(\mathcal H_0,\mathcal H\in\mathfrak B\),
\[
\mathcal H_0\cdot\mathcal H\in\mathfrak B.
\]
Lorsque \(\mathfrak B=\mathfrak M\), cette stabilité est automatique. C'est elle qui permet de transformer l'équivalence en réponse en congruence comportementale complète.

Les deux dernières propriétés sont distinctes. La fermeture par préfixes permet de remonter vers le passé d'une histoire admise ; la stabilité par continuation permet de prolonger un état atteint sans sortir de la classe déclarée.

\subsection{Classe des continuations après un préfixe}
\label{annexe-histoires-continuations-prefixe}

Pour \(\mathcal H_0\in\mathfrak M\), on définit la classe des
continuations admissibles après le préfixe \(\mathcal H_0\) par
\begin{equation}
\label{eq:annexe-G-quotient-gauche}
\mathcal H_0^{-1}\mathfrak B
=
\left\{
\mathcal H\in\mathfrak M
\;\middle|\;
\mathcal H_0\cdot\mathcal H\in\mathfrak B
\right\}.
\end{equation}
Il s'agit de la classe des continuations qui restent observables après
le préfixe \(\mathcal H_0\). Cette classe, et non \(\mathfrak B\)
elle-même, est en général le domaine pertinent lorsque l'on transporte
une comparaison vers un état déjà transformé. La notation
\(\mathcal H_0^{-1}\mathfrak B\) est conservée comme notation ensembliste ;
elle n'est pas employée ici pour qualifier latéralement la congruence.

\section{Exécutions, résolutions et lectures}
\label{annexe-histoires-observation}

Pour une histoire \(\mathcal H\), on note
\[
\operatorname{Exec}_{\mathcal H}(X)
\]
l'exécution complète issue de \(X\). Elle peut contenir la trajectoire
continue, les instants d'événement, les états pré- et post-saut et
l'état terminal. Le détail de cet objet dépend du modèle, mais il est
supposé déterminé par \(X\) et \(\mathcal H\).

La composition des transformations d'état est complétée par l'axiome de
concaténation des exécutions :
\begin{equation}
\label{eq:annexe-G-concatenation-executions}
\boxed{
\operatorname{Exec}_{\mathcal H_1\cdot\mathcal H_2}(X)
=
\operatorname{Exec}_{\mathcal H_1}(X)
\frown
\operatorname{Exec}_{\mathcal H_2}
\!\left(\Phi_{\mathcal H_1}(X)\right),
}
\end{equation}
où \(\frown\) désigne la concaténation chronologique de deux segments
d'exécution compatibles. L'exécution de \(\mathcal H_2\) depuis l'état
atteint après \(\mathcal H_1\) est donc le suffixe terminal de
l'exécution composée ; elle n'est pas, en général, identique à
l'exécution complète de \(\mathcal H_1\cdot\mathcal H_2\).

Une résolution \(\rho\) associe à une exécution une observation dans un
espace \(\mathcal Y_\rho\). Deux types de lectures doivent être
distingués.

\paragraph{Lecture terminale.}
Une lecture terminale dépend seulement de l'état terminal, par une
application
\[
\tau_\rho:\mathfrak X\longrightarrow\mathcal Y_\rho,
\qquad
\mathcal R_X^\rho(\mathcal H)
=
\tau_\rho\!\left(\Phi_{\mathcal H}(X)\right).
\]
Elle est supposée invariante par prolongation de la relaxation : dès que
le régime terminal déclaré est atteint, prolonger une phase de
relaxation sans nouvel événement pertinent ne modifie pas la lecture.
Cette condition rend notamment bien définie la lecture du bassin
terminal.

\paragraph{Lecture de chemin.}
Une lecture de chemin est une application
\[
\omega_\rho:\operatorname{Exec}\longrightarrow\mathcal Y_\rho,
\qquad
\mathcal R_X^\rho(\mathcal H)
=
\omega_\rho\!\left(\operatorname{Exec}_{\mathcal H}(X)\right).
\]
Pour le transport par préfixe, une telle lecture doit être
\emph{compatible au suffixe}. Cela signifie que la lecture utilisée
après un préfixe porte sur le segment terminal correspondant à la
continuation et non sur l'intégralité du chemin antérieur. Formellement,
pour tout segment initial \(E_0\) et tout segment terminal \(E\) admis,
la lecture ancrée sur ce segment satisfait
\begin{equation}
\label{eq:annexe-G-compatibilite-suffixe}
\omega_\rho^{\mathrm{fin}}(E_0\frown E)
=
\omega_\rho(E).
\end{equation}
Cette condition est satisfaite par les lectures effectivement utilisées
dans la suite, telles que le nombre de Kairos produits par la sonde
terminale. Elle ne l'est pas par une lecture intégrale comme le nombre
total de Kairos de toute l'exécution composée.

Les résultats de cette annexe utilisent les conditions de cohérence
suivantes.

\begin{enumerate}
\item[\textnormal{(O1)}]
  Pour tout \(X\) et toute histoire admissible \(\mathcal H\),
  l'exécution \(\operatorname{Exec}_{\mathcal H}(X)\) et sa lecture
  \(\mathcal R_X^\rho(\mathcal H)\) sont déterminées.
\item[\textnormal{(O2)}]
  La lecture est non invasive : elle compare les exécutions sans modifier
  la dynamique qu'elle lit.
\item[\textnormal{(O2$'$)}]
  Toute lecture terminale est invariante par prolongation admissible de
  la relaxation après l'atteinte du régime terminal déclaré.
\item[\textnormal{(O3)}]
  Si \(\rho_2\) raffine \(\rho_1\), il existe une application
  \(\pi_{21}:\mathcal Y_{\rho_2}\to\mathcal Y_{\rho_1}\) telle que la
  lecture grossière s'obtienne par projection de la lecture fine.
\item[\textnormal{(O4)}]
  Les projections de lecture associées à trois résolutions emboîtées
  sont compatibles : si \(\rho_3\) raffine \(\rho_2\), qui raffine
  \(\rho_1\), alors
  \[
  \pi_{31}=\pi_{21}\circ\pi_{32}.
  \]
\end{enumerate}

La résolution binaire utilisée dans plusieurs chapitres est terminale :
elle lit seulement le bassin atteint, par exemple repos ou
appropriation. Une résolution plus fine peut joindre à cette issue une
lecture de chemin compatible au suffixe, comme le nombre de Kairos de la
sonde ou une signature de son calendrier.

\section{Équivalence et quotient d'observation}
\label{annexe-histoires-equivalence}

Soient une classe admissible \(\mathfrak B\) et une résolution
\(\rho\). Deux systèmes \(X,Y\in\mathfrak X\) sont
\emph{observationnellement équivalents relativement à
\((\mathfrak B,\rho)\)} si leurs réponses coïncident sur toute la
classe :
\begin{equation}
\label{eq:annexe-G-equivalence}
X\sim_{\mathfrak B,\rho}Y
\quad\Longleftrightarrow\quad
\forall\mathcal H\in\mathfrak B,
\qquad
\mathcal R_X^\rho(\mathcal H)
=
\mathcal R_Y^\rho(\mathcal H).
\end{equation}

Cette relation est réflexive, symétrique et transitive, puisqu'elle est
l'égalité de deux fonctions sur un même domaine. Elle définit donc le
quotient d'observation
\[
\mathfrak X/{\sim_{\mathfrak B,\rho}}.
\]
La classe de \(X\) est notée
\[
[X]_{\mathfrak B,\rho}.
\]

Le qualificatif \emph{observationnel} est essentiel. Deux états
physiquement distincts peuvent appartenir à la même classe si aucune
histoire de \(\mathfrak B\), lue à la résolution \(\rho\), ne les
sépare. Inversement, une même paire peut être confondue à une
résolution et séparée à une autre.

\begin{encadre}[title={Vocabulaire}]
À budget fini, on parle de \emph{quotient d'observation}. Le terme
\emph{quotient de Nerode} est réservé à la classe complète des
continuations admissibles, pour laquelle l'équivalence devient une
congruence comportementale. Le système effectif du Tome III est la
famille ordonnée des quotients d'observation, non un quotient fini
isolé supposé définitif.
\end{encadre}

\section{Transport par préfixe}
\label{annexe-histoires-transport}

Le transport repose sur la composition des états pour les lectures
terminales et, pour les lectures de chemin, sur la concaténation des
exécutions jointe à la compatibilité au suffixe.

\resultstep{lem:G-transport-prefixe}
\paragraph{Lemme \theproposition{} --- transport par préfixe.}
Soit \(\rho\) une lecture terminale ou une lecture de chemin compatible
au suffixe. Pour tous \(X\in\mathfrak X\) et
\(\mathcal H_0,\mathcal H\in\mathfrak M\), la réponse à la
continuation \(\mathcal H\) depuis l'état atteint après
\(\mathcal H_0\) coïncide avec la lecture du segment terminal
correspondant dans l'exécution composée. On note cette identité
\begin{equation}
\label{eq:annexe-G-identite-transport}
\mathcal R_{\Phi_{\mathcal H_0}(X)}^\rho(\mathcal H)
=
\mathcal R_X^{\rho,\mathrm{fin}}
  (\mathcal H_0\cdot\mathcal H\,;\mathcal H),
\end{equation}
où le membre de droite désigne la lecture ancrée sur le suffixe produit
par \(\mathcal H\). Dans le cas terminal, cette annotation est inutile
et l'identité se réduit à
\[
\tau_\rho\!\left(
\Phi_{\mathcal H}(\Phi_{\mathcal H_0}(X))
\right)
=
\tau_\rho\!\left(
\Phi_{\mathcal H_0\cdot\mathcal H}(X)
\right).
\]

Par conséquent,
\begin{equation}
\label{eq:annexe-G-transport}
X\sim_{\mathfrak B,\rho}Y
\quad\Longrightarrow\quad
\Phi_{\mathcal H_0}(X)
\sim_{\mathcal H_0^{-1}\mathfrak B,\rho}
\Phi_{\mathcal H_0}(Y),
\end{equation}
à condition que, lorsqu'il s'agit d'une lecture de chemin,
l'équivalence sur \(\mathfrak B\) soit définie avec la même convention
d'ancrage terminal pour les histoires préfixées.

\paragraph{Démonstration.}
Pour une lecture terminale, la convention de composition donne
\[
\Phi_{\mathcal H}(\Phi_{\mathcal H_0}(X))
=
\Phi_{\mathcal H_0\cdot\mathcal H}(X),
\]
et l'application de \(\tau_\rho\) établit immédiatement
\eqref{eq:annexe-G-identite-transport}.

Pour une lecture de chemin, l'axiome
\eqref{eq:annexe-G-concatenation-executions} donne
\[
\operatorname{Exec}_{\mathcal H_0\cdot\mathcal H}(X)
=
\operatorname{Exec}_{\mathcal H_0}(X)
\frown
\operatorname{Exec}_{\mathcal H}
\!\left(\Phi_{\mathcal H_0}(X)\right).
\]
La compatibilité au suffixe
\eqref{eq:annexe-G-compatibilite-suffixe} identifie alors la lecture du
segment terminal à celle de l'exécution de \(\mathcal H\) depuis
\(\Phi_{\mathcal H_0}(X)\), ce qui établit la même identité.

Supposons maintenant \(X\sim_{\mathfrak B,\rho}Y\), et prenons
\(\mathcal H\in\mathcal H_0^{-1}\mathfrak B\). Par définition,
\(\mathcal H_0\cdot\mathcal H\in\mathfrak B\). L'égalité des réponses
sur cette histoire, avec la convention terminale déclarée, puis
l'identité précédente appliquée à \(X\) et à \(Y\), donnent
\[
\mathcal R_{\Phi_{\mathcal H_0}(X)}^\rho(\mathcal H)
=
\mathcal R_{\Phi_{\mathcal H_0}(Y)}^\rho(\mathcal H).
\]
Cette égalité valant pour toute continuation de
\(\mathcal H_0^{-1}\mathfrak B\),
\eqref{eq:annexe-G-transport} est démontrée.
\hfill\(\square\)

\paragraph{Portée.}
L'identité ne vaut pas pour une lecture de chemin intégrale qui conserve
le préfixe dans la quantité observée. Par exemple, le nombre total de
Kairos de l'exécution \(\mathcal H_0\cdot\mathcal H\) peut différer du
nombre de Kairos produits par la seule continuation \(\mathcal H\).
Cette restriction ne touche ni les lectures terminales ni les lectures
de sonde employées dans le Tome III.

Une équivalence à budget fini n'est pas, en général, conservée sur le
même budget après l'application d'un préfixe. Elle est transportée vers
la classe exacte des continuations encore testables. Pour obtenir
\[
\Phi_{\mathcal H_0}(X)
\sim_{\mathfrak B,\rho}
\Phi_{\mathcal H_0}(Y),
\]
il suffit naturellement que
\[
\mathfrak B\subseteq\mathcal H_0^{-1}\mathfrak B,
\]
c'est-à-dire que toute histoire de \(\mathfrak B\) reste admissible
après \(\mathcal H_0\). Cette condition suffisante n'est pas satisfaite
par la plupart des budgets finis ; l'implication peut néanmoins tenir
pour certaines paires par coïncidence de leurs réponses.

\section{Congruence comportementale complète}
\label{annexe-histoires-nerode}

On considère maintenant le monoïde libre complet
\(\mathfrak M=\mathcal E^\ast\). Sa stabilité par concaténation découle
de sa définition.

\resultstep{prop:G-congruence-nerode}
\paragraph{Proposition \theproposition{} --- congruence de Nerode.}
Soit \(\rho\) une lecture terminale ou une lecture de chemin compatible
au suffixe. La relation
\[
X\sim_{\mathfrak M,\rho}Y
\quad\Longleftrightarrow\quad
\forall\mathcal H\in\mathfrak M,
\quad
\mathcal R_X^\rho(\mathcal H)
=
\mathcal R_Y^\rho(\mathcal H)
\]
est une congruence pour l'action des histoires : pour tout
\(\mathcal H_0\in\mathfrak M\),
\begin{equation}
\label{eq:annexe-G-nerode-congruence}
X\sim_{\mathfrak M,\rho}Y
\quad\Longrightarrow\quad
\Phi_{\mathcal H_0}(X)
\sim_{\mathfrak M,\rho}
\Phi_{\mathcal H_0}(Y).
\end{equation}

\paragraph{Démonstration.}
Puisque \(\mathfrak M\) est le monoïde libre complet,
\[
\mathcal H_0^{-1}\mathfrak M=\mathfrak M.
\]
Le lemme de transport par préfixe
\eqref{eq:annexe-G-transport} donne alors directement
\eqref{eq:annexe-G-nerode-congruence}.
\hfill\(\square\)

Le quotient
\[
\mathfrak X/{\sim_{\mathfrak M,\rho}}
\]
est le quotient comportemental complet à la résolution \(\rho\). Deux
classes distinctes possèdent nécessairement une histoire séparatrice ;
deux représentants d'une même classe restent indiscernables après
toute histoire préalable et sous toute continuation admissible.

La disponibilité de toutes les continuations est utilisée exactement
ici. Les autres résultats fondamentaux de cette annexe restent valides
sur une classe minimale.

\section{Budgets et graduations des histoires}
\label{annexe-histoires-budgets}

Une classe \(\mathfrak B\) peut être construite à partir d'une mesure de
complexité, mais elle ne s'y réduit pas. Parmi les graduations utilisées
ou envisageables figurent :
\[
|\mathcal H|_{\mathrm{ep}}
=
\text{nombre d'épisodes},
\]
\[
|\mathcal H|_{\mathrm{dur}}
=
\text{durée active},
\qquad
|\mathcal H|_{\mathrm{tot}}
=
\text{durée totale},
\]
et une dose cumulée telle que
\[
|\mathcal H|_{\mathrm{dose}}
=
\sum_k g_kT_k.
\]

Il est souvent préférable de conserver une complexité vectorielle
\begin{equation}
\label{eq:annexe-G-complexite-vectorielle}
\mathbf c(\mathcal H)
=
\bigl(
 n_{\mathrm{ep}},
 T_{\mathrm{actif}},
 T_{\mathrm{total}},
 Q
\bigr)
\end{equation}
et de définir un budget par un ensemble inférieur
\begin{equation}
\label{eq:annexe-G-budget-vectoriel}
\mathfrak B_{\mathbf b}
=
\left\{
\mathcal H\in\mathfrak M
\;\middle|\;
\mathbf c(\mathcal H)\preceq\mathbf b
\right\}.
\end{equation}
Cette construction évite d'imposer trop tôt une compensation arbitraire
entre le nombre d'épisodes, leur durée et leur intensité.

Une fonction scalaire
\[
c:\mathfrak M\longrightarrow\mathbb R_+
\]
peut néanmoins être introduite pour ordonner les histoires séparatrices
ou construire une distance. Elle ne participe pas à la définition de
l'équivalence sur une classe déjà déclarée.

\section{Indépendance du quotient vis-à-vis de la complexité}
\label{annexe-histoires-c-independance}

\resultstep{prop:G-c-independance-quotient}
\paragraph{Proposition \theproposition{} --- \(c\)-indépendance du quotient.}
Pour une classe \(\mathfrak B\) et une résolution \(\rho\) fixées, la
relation \(\sim_{\mathfrak B,\rho}\) ne dépend d'aucune fonction de
complexité \(c\). En revanche, la profondeur à laquelle une paire est
séparée dépend de \(c\).

Cette indépendance est une propriété à classe fixée. Si le budget est
lui-même défini par la complexité, par exemple
\(
\mathfrak B_b^c=\{\mathcal H:c(\mathcal H)\leq b\}
\), changer \(c\) peut changer \(\mathfrak B_b^c\), donc le quotient
d'observation. Ce changement passe par la classe admissible et non par
la définition de l'équivalence sur une classe donnée.

\paragraph{Démonstration.}
Par définition,
\[
X\sim_{\mathfrak B,\rho}Y
\]
si et seulement si les réponses de \(X\) et \(Y\) coïncident pour
toutes les histoires de \(\mathfrak B\). La fonction \(c\) n'apparaît
pas dans cette condition. Elle ne peut donc modifier ni les paires
équivalentes ni le quotient associé.

Pour une paire non équivalente, on peut en revanche définir
\begin{equation}
\label{eq:annexe-G-longueur-separation}
\ell_{c,\mathfrak B,\rho}(X,Y)
=
\inf\left\{
 c(\mathcal H)
 \;\middle|\;
 \mathcal H\in\mathfrak B,
 \ \mathcal R_X^\rho(\mathcal H)
 \neq
 \mathcal R_Y^\rho(\mathcal H)
\right\}.
\end{equation}
Cette valeur change lorsque \(c\) change, même si l'ensemble des
histoires séparatrices reste identique. La séparation répond donc à la
question \emph{les systèmes diffèrent-ils dans ce protocole ?}, tandis
que \(\ell\) répond à la question \emph{à quelle profondeur la mesure
choisie rend-elle cette différence visible ?}.
\hfill\(\square\)

Si aucune histoire de \(\mathfrak B\) ne sépare la paire, on adopte la
convention
\[
\ell_{c,\mathfrak B,\rho}(X,Y)=+\infty.
\]
La quantité \(\ell\) est ainsi toujours relative au quadruplet
\((X,Y,\mathfrak B,\rho)\), puis à la graduation \(c\). Elle n'est pas
un invariant absolu du couple de systèmes.

\section{Monotonie bidimensionnelle de l'observabilité}
\label{annexe-histoires-monotonie}

On ordonne les budgets par inclusion. On dit que \(\rho_2\) raffine
\(\rho_1\), et l'on note
\[
\rho_1\preceq\rho_2,
\]
lorsque la condition (O3) est satisfaite.

\resultstep{prop:G-monotonie-observabilite}
\paragraph{Proposition \theproposition{} --- monotonie bidimensionnelle.}
Si
\[
\mathfrak B_1\subseteq\mathfrak B_2
\qquad\text{et}\qquad
\rho_1\preceq\rho_2,
\]
alors
\begin{equation}
\label{eq:annexe-G-monotonie}
\boxed{
\sim_{\mathfrak B_2,\rho_2}
\ \subseteq\ 
\sim_{\mathfrak B_1,\rho_1}.
}
\end{equation}
Toute augmentation de ce que l'on peut faire vivre ou de ce que l'on
peut lire peut scinder des classes d'équivalence, mais ne peut pas en
fusionner.

\paragraph{Démonstration.}
Soient \(X,Y\in\mathfrak X\) tels que
\(X\sim_{\mathfrak B_2,\rho_2}Y\). Pour toute
\(\mathcal H\in\mathfrak B_1\), l'inclusion des budgets donne
\(\mathcal H\in\mathfrak B_2\), donc
\[
\mathcal R_X^{\rho_2}(\mathcal H)
=
\mathcal R_Y^{\rho_2}(\mathcal H).
\]
Comme \(\rho_2\) raffine \(\rho_1\), il existe \(\pi_{21}\) telle que
\[
\mathcal R_Z^{\rho_1}(\mathcal H)
=
\pi_{21}\!\left(
\mathcal R_Z^{\rho_2}(\mathcal H)
\right)
\]
pour tout système \(Z\). En appliquant \(\pi_{21}\) à l'égalité
précédente, on obtient
\[
\mathcal R_X^{\rho_1}(\mathcal H)
=
\mathcal R_Y^{\rho_1}(\mathcal H).
\]
Cette égalité valant pour toute histoire de \(\mathfrak B_1\), on a
\(X\sim_{\mathfrak B_1,\rho_1}Y\).
\hfill\(\square\)

Deux cas particuliers sont utiles :
\[
\mathfrak B_1\subseteq\mathfrak B_2
\quad\Longrightarrow\quad
\sim_{\mathfrak B_2,\rho}
\subseteq
\sim_{\mathfrak B_1,\rho},
\]
et
\[
\rho_1\preceq\rho_2
\quad\Longrightarrow\quad
\sim_{\mathfrak B,\rho_2}
\subseteq
\sim_{\mathfrak B,\rho_1}.
\]
Le premier raffinement agit en sollicitant davantage le système ; le
second en lisant plus finement une même sollicitation.

\section{Système inverse des quotients d'observation}
\label{annexe-histoires-systeme-inverse}

Considérons l'ensemble ordonné des couples \((\mathfrak B,\rho)\), avec
\begin{equation}
\label{eq:annexe-G-ordre-produit}
(\mathfrak B_1,\rho_1)
\preceq
(\mathfrak B_2,\rho_2)
\quad\Longleftrightarrow\quad
\mathfrak B_1\subseteq\mathfrak B_2
\ \text{et}\ 
\rho_1\preceq\rho_2.
\end{equation}
La proposition~\ref{prop:G-monotonie-observabilite} fournit, pour chaque paire ordonnée, une
application naturelle
\begin{equation}
\label{eq:annexe-G-projection-quotients}
q_{21}:
\mathfrak X/{\sim_{\mathfrak B_2,\rho_2}}
\longrightarrow
\mathfrak X/{\sim_{\mathfrak B_1,\rho_1}},
\qquad
q_{21}\!\left([X]_{\mathfrak B_2,\rho_2}\right)
=
[X]_{\mathfrak B_1,\rho_1}.
\end{equation}
Cette application est bien définie grâce à
\eqref{eq:annexe-G-monotonie}. Pour trois niveaux emboîtés, la
définition même de \(q_{ji}([X])=[X]\) donne automatiquement
\[
q_{31}=q_{21}\circ q_{32}.
\]
La condition (O4) assure, en amont, que les projections de lecture
\(\pi_{ji}\) utilisées pour comparer les résolutions sont elles-mêmes
cohérentes. Les quotients forment donc un système inverse.

Un élément de la limite projective, lorsqu'on considère formellement
tous les niveaux, est une famille cohérente
\[
\left(C_{\mathfrak B,\rho}\right)_{(\mathfrak B,\rho)}
\]
de classes telle que chaque projection d'une classe fine redonne la
classe plus grossière correspondante. Cette limite ne suppose pas que
la filtration se stabilise à un budget fini.

\begin{encadre}[title={Lecture effective}]
L'objet de travail du Tome III est le système inverse des quotients
d'observation. Un quotient à budget fixé en est une coupe. Le quotient
comportemental complet est une limite idéale lorsque toutes les
continuations et la résolution déclarée sont prises en compte.
\end{encadre}

Cette formulation évite deux confusions. D'une part, l'absence de
nouvelle séparation sur quelques budgets successifs ne prouve pas la
stationnarité définitive. D'autre part, une classe observée à un niveau
ne doit pas être essentialisée : elle peut se scinder dès que le
protocole ou la lecture sont enrichis.

\section{Conséquences pour les résumés et les seuils}
\label{annexe-histoires-consequences}

Soit
\[
Z_{\mathfrak B,\rho}:\mathfrak X\longrightarrow\mathcal Z
\]
un résumé déclaré suffisant à un niveau d'observation. Cette
suffisance signifie
\begin{equation}
\label{eq:annexe-G-resume-suffisant}
Z_{\mathfrak B,\rho}(X)
=
Z_{\mathfrak B,\rho}(Y)
\quad\Longrightarrow\quad
X\sim_{\mathfrak B,\rho}Y.
\end{equation}
Elle ne vaut pas automatiquement après enrichissement du budget ou de
la résolution. La proposition~\ref{prop:G-monotonie-observabilite} autorise seulement l'implication
dans le sens grossier : un résumé suffisant à un niveau fin reste
suffisant après projection vers un niveau moins exigeant.

De même, un ensemble de commutation \(\Theta\), une surface de frontière
ou une horloge déclarée active doivent porter les indices
\((\mathfrak B,\rho)\) dès que leur définition repose sur les réponses
observées. Deux résolutions peuvent faire apparaître des séparations à
des niveaux d'histoire différents. Deux budgets peuvent donner des
valeurs différentes de la profondeur de séparabilité
\(\ell_{c,\mathfrak B,\rho}\), toutes deux correctes dans leur domaine.

Cette dépendance n'est pas une faiblesse du formalisme. Elle constitue
son contenu : l'individualité comportementale est ce qui résiste à une
famille déclarée d'interventions et de lectures.

\section{Niveaux de classe requis par les résultats fondamentaux}
\label{annexe-histoires-table-portee}

Le tableau suivant résume les hypothèses de fermeture réellement
utilisées.

\begin{center}
\begin{longtable}{p{0.24\linewidth}p{0.25\linewidth}p{0.39\linewidth}}
\toprule
Résultat & Niveau requis & Raison \\
\midrule
\endfirsthead
\toprule
Résultat & Niveau requis & Raison \\
\midrule
\endhead
Définition de
\(\sim_{\mathfrak B,\rho}\)
& Classe minimale
& Il suffit que les réponses soient définies sur
  \(\mathfrak B\). \\
\addlinespace
Lemme~\ref{lem:G-transport-prefixe}
& Classe minimale
& Le résultat transporte l'équivalence vers la classe de continuations
  \(\mathcal H_0^{-1}\mathfrak B\), sans supposer qu'il coïncide avec
  \(\mathfrak B\). \\
\addlinespace
Proposition~\ref{prop:G-congruence-nerode}
& Stabilité par continuation
& Elle assure
  \(\mathcal H_0^{-1}\mathfrak M=\mathfrak M\) et transforme le
  transport par préfixe en congruence sur la même classe. \\
\addlinespace
Proposition~\ref{prop:G-monotonie-observabilite}
& Classes minimales ordonnées
& Seules l'inclusion des budgets et la projection des résolutions sont
  utilisées. \\
\addlinespace
Proposition~\ref{prop:G-c-independance-quotient}
& Classe minimale
& La preuve repose uniquement sur la définition de l'égalité des
  réponses. \\
\bottomrule
\end{longtable}
\end{center}

Ainsi, parmi ces résultats, seule la proposition de congruence complète
requiert la stabilité par continuation. Les mécanismes de seuil, les
réductions à une sonde et les frontières construites dans les chapitres
suivants peuvent être établis dans les classes minimales explicitement
déclarées.

\section{Table des notations}
\label{annexe-histoires-notations}

\begin{center}
\begin{longtable}{p{0.23\linewidth}p{0.65\linewidth}}
\toprule
Notation & Signification \\
\midrule
\endfirsthead
\toprule
Notation & Signification \\
\midrule
\endhead
\(\mathfrak X\)
& Ensemble des états complets ou des présentations de systèmes. \\
\addlinespace
\(\mathcal E\)
& Alphabet des épisodes élémentaires, comprenant épisodes actifs et délais. \\
\addlinespace
\(\mathcal E_{\mathrm{act}}\)
& Alphabet des épisodes actifs. \\
\addlinespace
\(\mathcal D\)
& Alphabet des délais, considérés comme générateurs à part entière. \\
\addlinespace
\(\mathfrak M=\mathcal E^\ast\)
& Monoïde des histoires finies. \\
\addlinespace
\(\varepsilon\)
& Histoire vide. \\
\addlinespace
\(\mathcal H_1\cdot\mathcal H_2\)
& Concaténation chronologique : \(\mathcal H_1\) puis
  \(\mathcal H_2\). \\
\addlinespace
\(\Phi_{\mathcal H}\)
& Transformation de l'état produite par l'histoire \(\mathcal H\). \\
\addlinespace
\(\mathfrak B\)
& Classe ou budget d'histoires admissibles. \\
\addlinespace
\(\mathcal H_0^{-1}\mathfrak B\)
& Continuations \(\mathcal H\) telles que
  \(\mathcal H_0\cdot\mathcal H\in\mathfrak B\). \\
\addlinespace
\(\rho\)
& Résolution d'observation. \\
\addlinespace
\(\omega_\rho\)
& Lecture d'une exécution à la résolution \(\rho\). \\
\addlinespace
\(\mathcal R_X^\rho(\mathcal H)\)
& Réponse du système \(X\) à l'histoire \(\mathcal H\). \\
\addlinespace
\(X\sim_{\mathfrak B,\rho}Y\)
& Égalité des réponses de \(X\) et \(Y\) sur tout
  \(\mathfrak B\), à la résolution \(\rho\). \\
\addlinespace
\([X]_{\mathfrak B,\rho}\)
& Classe d'équivalence observationnelle de \(X\). \\
\addlinespace
\(c,\mathbf c\)
& Complexité scalaire ou vectorielle d'une histoire. \\
\addlinespace
\(\ell_{c,\mathfrak B,\rho}(X,Y)\)
& Profondeur de séparabilité de \(X\) et \(Y\) : infimum des complexités
  des histoires séparatrices, ou \(+\infty\) si aucune histoire de
  \(\mathfrak B\) ne les sépare. \\
\addlinespace
\(q_{21}\)
& Projection naturelle d'un quotient fin vers un quotient plus
  grossier. \\
\bottomrule
\end{longtable}
\end{center}

\section{Conclusion de l'annexe}
\label{annexe-histoires-conclusion}

Le cadre peut être résumé par les implications suivantes :
\[
\text{composition des histoires et compatibilité de la lecture}
\Longrightarrow
\text{transport par préfixe},
\]
\[
\text{stabilité par continuation}
\Longrightarrow
\text{congruence de Nerode},
\]
\[
\text{budget plus riche ou résolution plus fine}
\Longrightarrow
\text{quotient plus fin},
\]
\[
\text{mesure de complexité}
\Longrightarrow
\text{géométrie de la séparation, non définition du quotient}.
\]

La suite du Tome III peut dès lors s'appuyer sur un langage unique. Une
histoire n'est jamais séparatrice en soi, mais relativement à une
résolution. Une classe n'est jamais définitive, mais située dans une
filtration. Un résumé n'est jamais suffisant absolument, mais pour un
couple \((\mathfrak B,\rho)\) déclaré. Ces précautions formelles sont
ce qui permet ensuite de parler rigoureusement de contenu survivant,
d'état lent, de frontière d'appropriation et de plasticité.
